\textbf{Кому подходит схема.} Расчёт относится к команде с уже имеющимися M5 Max 64~ГБ,
собственным выделенным RTX PRO 6000 с vLLM FP8 и предполагаемым каналом 125~МБ/с, RTT 1~мс.
Лёгкая нагрузка составляет 200 вызовов в день, средняя 800; тяжёлая, 2500, соответствует интенсивности владельца.
Лёгкая и средняя нагрузки типичны. Доля открытой модели 20\%; сравниваются формы вызовов внутри общей измеренной лестницы,
10\,885 из 39\,904 (27.28\%), с приведением к заданному числу вызовов. Полная исходная смесь этим расчётом не оценена.

Для команды, которой нельзя передавать данные наружу, гибрид дешевле своего кластера на том же сервере
при средней нагрузке с пяти человек, при тяжёлой с двух, в проверенном диапазоне до 500 человек.
При средней нагрузке и пяти людях цена \$130.33 на человека в месяц при рабочей/календарной занятости
6.19/1.42\% против \$245.62 и 32.17/7.40\% у своего кластера. Ответы генерируют Маки, поэтому серверов требуется меньше.
При тяжёлой нагрузке собственный офисный сервер даёт выигрыш против медианного допущенного API
с 45 человек при условной доле кэша 98.1\%; на арендованном сервере с каналом 112.5~МБ/с и RTT 400~мс
такого пересечения до 500 человек нет. Для общего занятого сервера предел распределённых расходов
при доле кэша 98.1\% даёт отношение цены медианного API к гибриду 1.701 при тяжёлой нагрузке
и 0.934 при средней; эти значения используют распределение при занятости 70\%, нижняя оценка при 100\% приведена отдельно.
Сервис общего prefill ведут и авторы статьи; выделенный сервер остаётся основной базой сравнения.

Лёгкая нагрузка оставляет выделенный сервер команды из 50 человек занятым на 3.61\% рабочего времени
и 0.83\% календарного; медианный API при условных 98.1\% кэша примерно в 11 раз дешевле.
Свой кластер на том же сервере также дешевле: обоим маршрутам хватает одинакового парка, гибрид добавляет расходы клиента.
Самый дешёвый допущенный API при условной доле кэша 98.1\% дешевле гибрида примерно в 55 раз при лёгкой нагрузке
и в 4.4 раза при тяжёлой. Без кэша гибрид обходит самый дешёвый допущенный API только при тяжёлой нагрузке
с 74 человек на офисном сервере. Реальная доля кэша поставщиков ещё измеряется.
Покупка M5 Max добавляет \$69.05 на человека в месяц: пересечение со своим кластером при средней нагрузке
сдвигается с 5 до 34 человек, пересечение с медианным API при условных 98.1\% кэша исчезает.

\textbf{Ожидание и автономность.} Выигрыш M3 Max в статье измерен во времени; строк экономики M3 Max здесь нет.
Порция 10 тысяч токенов на глубине 150 тысяч при обычном приоритете требует около $10000/44.4=225$ секунд.
Передача в bf16 несёт $10000\times34816+78446592=426606592$ байта, около 426.6~МБ:
при 125~МБ/с только передача занимает $426.6/125=3.4$ секунды, при 5~МБ/с $426.6/5=85.3$;
к передаче добавляются серверный prefill и импорт. M5 Max читает быстрее, поэтому на 100 тысячах и порции 2343
с bf16 порог канала равен 34.475~МБ/с против 5.305 у M3 Max при RTT 1~мс.
Измеренный предварительный расчёт на канале 5.19~МБ/с с RTT 125.8~мс оставил все ходы локальными.
На миллионе токенов оценка порога M3 Max с bf16 равна 2.992~МБ/с; прямой замер ещё ожидается.
Без сервера Мак продолжает работу медленнее; серверная копия сессии позволяет восстановиться после сбоя.

Для владельца ПК с RTX 5090 гибрид мало добавляет при модели 27B и контексте до 262 тысяч:
одна сессия помещается в 32~ГБ, prefill быстрее M3 Max примерно в 23 раза, генерация в 4.3--4.7 раза.
Граница применимости проходит там, где модель и контекст не помещаются в видеопамять: миллион токенов,
несколько сессий одновременно; для крупных MoE это ожидание без замера.
Чужой сервер prefill видит весь контекст пользователя. Данные остаются внутри компании при собственном сервере.
